% Acronym definitions using the `acro` package.
% Add `\usepackage{acro}` to your preamble (main.tex) if not already present.

\DeclareAcronym{vsb}{
	short = {VŠB},
	long  = {Vysoká škola Báňská}
}

\DeclareAcronym{vcs}{
	short = {VCS},
	long  = {Version Control System}
}

\DeclareAcronym{api}{
	short = {API},
	long  = {Application Programming Interface}
}

\DeclareAcronym{cli}{
	short = {CLI},
	long  = {Command Line Interface}
}

\DeclareAcronym{gui}{
	short = {GUI},
	long  = {Graphical User Interface}
}

\DeclareAcronym{ide}{
	short = {IDE},
	long  = {Integrated Development Environment}
}

\DeclareAcronym{http}{
	short = {HTTP},
	long  = {Hypertext Transfer Protocol}
}

\DeclareAcronym{ssh}{
	short = {SSH},
	long  = {Secure Shell}
}

\DeclareAcronym{url}{
	short = {URL},
	long  = {Uniform Resource Locator}
}

\DeclareAcronym{json}{
	short = {JSON},
	long  = {JavaScript Object Notation}
}

\DeclareAcronym{git}{
	short = {Git},
	long  = {Git}
}

\DeclareAcronym{svn}{
	short = {SVN},
	long  = {Subversion}
}

\DeclareAcronym{cvs}{
	short = {CVS},
	long  = {Concurrent Versions System}
}

\DeclareAcronym{sha}{
	short = {SHA},
	long  = {Secure Hash Algorithm}
}

\DeclareAcronym{tcpip}{
	short = {TCP/IP},
	long  = {Transmission Control Protocol/Internet Protocol}
}

\DeclareAcronym{https}{
	short = {HTTPS},
	long  = {Hypertext Transfer Protocol Secure}
}

% Shortcut macros: use `\API`, `\VSB`, etc. in your document.
\newcommand{\VSB}{\ac{vsb}}
\newcommand{\VCS}{\ac{vcs}}
\newcommand{\API}{\ac{api}}
\newcommand{\CLI}{\ac{cli}}
\newcommand{\GUI}{\ac{gui}}
\newcommand{\IDE}{\ac{ide}}
\newcommand{\HTTP}{\ac{http}}
\newcommand{\SSH}{\ac{ssh}}
\newcommand{\URL}{\ac{url}}
\newcommand{\JSON}{\ac{json}}
\newcommand{\GIT}{\ac{git}}
\newcommand{\SVN}{\ac{svn}}
\newcommand{\CVS}{\ac{cvs}}
\newcommand{\SHA}{\ac{sha}}
\newcommand{\TCPIP}{\ac{tcpip}}
\newcommand{\HTTPS}{\ac{https}}

